%%%%%%%%%%%%%%%%%%%%%%%%%%%%%%%%%%%%%%%%%%%%%%%%%%%%%%%%%%%%%%%%%%%%%%%%%%%
% Englische Fassung des Uebungsblatts zu Lehreinheit 7 (Teil A und B).
% Vorlage: handouts/LE07.tex. Nur in handouts.tex mit \def\wuzsprache{en}.
% Noten im Fallbeispiel: deutsche Notenskala (1 = beste, 5 = nicht
% bestanden), im Text erklaert.
%%%%%%%%%%%%%%%%%%%%%%%%%%%%%%%%%%%%%%%%%%%%%%%%%%%%%%%%%%%%%%%%%%%%%%%%%%%
\begin{uebungsblatt}{Research Design}

\aufgabe{Case study: analysing a design}[in class, pairs, 15 min.]
\textit{The following study is fictitious and was written for this
exercise; all numbers are invented.}

\begin{center}
\fbox{\parbox{0.92\textwidth}{\small
\textbf{\enquote{Learning app improves maths grades by 0.7}}\\[1mm]
In the winter semester, a university offers a new practice app for the
lecture Mathematics~1. The app was developed by the lecturer himself. Use
is voluntary; 64 of 212 students sign up. To introduce the app, the
lecturer offers the registered students a weekly tutorial in which they
practise with the app. At the end of the semester, all 212 students take
the same exam, which the lecturer marks himself; name and student number
are on each script. The app group achieves an average grade of 2.3, the
other 148 students 3.0 (German scale: 1 = best, 5 = fail). The
university's press release says: \enquote{The app improves the grade by
0.7. From next semester it will be compulsory for all first-year
students.}
}}
\end{center}

\begin{enumerate}
  \item Name the IV and DV (independent and dependent variable). How was
    the DV measured?
  \item What type of design is this? Write the design in R/X/O notation
    (one line per group; R = random assignment, X = treatment,
    O = observation).
  \item Is there a control group? Apart from using the app, in which
    respects does it differ from the app group?
  \item Name at least three confounders or threats to internal validity
    and match them to technical terms (e.g.\ selection, history,
    maturation, testing, instrumentation, attrition, researcher
    expectations).
  \item Is the measurement objective? What would a blind design look like?
  \item Sketch an improved design that could test the claim in the press
    release.
\end{enumerate}
\schreiblinien{9}

\aufgabe{Design sketch for your hypothesis}[start in class, due by Unit~8, graded]
Fill in the template and justify your choices on at most one further page.

\begin{center}
\renewcommand{\arraystretch}{1.5}
\begin{tabular}{|p{0.3\textwidth}|p{0.6\textwidth}|}
\hline
H\textsubscript{1} (from Unit~5) & \\[6mm] \hline
IV and its levels & \\ \hline
DV and measuring instrument & \\ \hline
Type of design (experimental / quasi / non-experimental; cross-sectional or longitudinal) & \\ \hline
Design in notation (R / X / O) & \\ \hline
Sample: who, how many, how recruited? & \\ \hline
Analysis plan (from Unit~6): test or Bayes factor, smallest effect of interest, $n$ from the power table & \\[6mm] \hline
Assignment to conditions & \\ \hline
Control condition (active?) & \\ \hline
Main confounders (dragons on your map from Unit~6) and how I control them & \\[8mm] \hline
Measurement quality: objective? reliable? valid? blind? parallel measurement? & \\[8mm] \hline
Remaining threat that I accept & \\ \hline
\end{tabular}
\end{center}
\end{uebungsblatt}

\begin{uebungsblattb}{3}{Feedback and Transfer}

\aufgabe{Peer feedback}[in class, pairs, 10 min.]
For your partner's sketch, note:
\begin{center}
\renewcommand{\arraystretch}{2}
\begin{tabular}{|p{0.42\textwidth}|p{0.48\textwidth}|}
\hline
One alternative explanation the design does not rule out & \\ \hline
One change that would rule it out & \\ \hline
What is convincing about the design & \\ \hline
\end{tabular}
\end{center}

\aufgabe{Diary of wonder}[homework]
Find a claim about an effect in advertising, the news or social media
that lacks a comparison. Which control condition would you need to test
it?

\begin{lehrhinweis}[Expected answers]
\textbf{Task~1.} IV: app use (yes/no, self-selected); DV: exam grade,
marked by the app's developer, not anonymised. Design: quasi-experimental
or non-experimental group comparison without pretest
(X\,O / \phantom{X\,}O), no randomisation. The \enquote{control group}
differs in motivation and probably prior knowledge (self-selection), in
the additional tutorial (the IV is confounded with practice time and
support), and in how the marker perceives them (experimenter effect, lack
of objectivity). Further points: no pretest, so no statement about change;
possible novelty effects; generalising from volunteers to compulsory use
is questionable (external validity). Improved design: random assignment
(possibly stratified by prior knowledge from an entry test), a tutorial of
equal length with conventional exercises for both groups (active
control), anonymised exam, marking by someone who does not know the group
membership (blind design).
\textbf{Task~2.} Assessed: fit between design and hypothesis (a causal
claim requires an experimental or a well-justified quasi-experimental
design), a control condition that differs only in the IV, a concrete
measuring instrument, at least two named and controlled confounders, an
honestly named remaining threat.
\end{lehrhinweis}
\end{uebungsblattb}
